\documentclass[UTF8,11pt]{ctexart}
\usepackage[a4paper,margin=22mm]{geometry}
\usepackage{graphicx,enumitem,amssymb,tabularx}
\usepackage[normalem]{ulem}
\usepackage{xeCJKfntef}
\setlength{\parindent}{0pt}
\setlength{\parskip}{.6em}
\sloppy
\begin{document}
\section*{2018年6月英语四级真题(第2套)}
\subsection*{\textbf{2018} 年\textbf{6} 月大学英语四级真题（第\textbf{2} 套\textbf{)}}


\subsection*{\textbf{Part I} \textbf{Writing} （\textbf{30 minutes}）}


\textbf{Directions:}\textit{ For this part, you are allowed 30minutes to write a short essay on}\textbf{\textit{ the importance}}


\textbf{\textit{of writing ability and how to develop it.}}\textit{ You should write at least 120 words but no more than} \textit{180 words.} \par\noindent\rule{.95\linewidth}{.4pt}\par  \par\noindent\rule{.95\linewidth}{.4pt}\par  \par\noindent\rule{.95\linewidth}{.4pt}\par


\subsection*{\textbf{Part II} \textbf{Listening Comprehension} \textbf{(25 minutes)}}


\subsection*{\textbf{Section A}}


\textbf{Directions:}\textit{ In this section, you will hear three news reports. At the end of each news report, you}


\textit{will hear two or three questions. Both the news report and the questions will be spoken only} \textit{once. After you hear a question, you must choose the best answer from the four choices} \textit{marked A), B), C) and D). Then mark the corresponding letter on}\textbf{\textit{ Answer Sheet 1}}\textit{ with a} \textit{single line through the centre.}


\textbf{Questions 1 and 2 are based on the news report you have just heard.}


\noindent\begin{tabularx}{\linewidth}{@{}lXXXX@{}}1. & \textbf{A.} Annoyed. & \textbf{B.} Scared. & \textbf{C.} Confused. & \textbf{D.} Offended.\\\end{tabularx}\par

2.


\begin{itemize}[leftmargin=2em,itemsep=.2em,topsep=.4em]\item[A.] It crawled over the woman's hands.
\item[B.] It wound up on the steering wheel.
\item[C.] It was killed by the police on the spot.
\item[D.] It was covered with large scales.
\end{itemize}

\textbf{Questions 3 and 4 are based on the news report you have just heard.}


3.


\begin{itemize}[leftmargin=2em,itemsep=.2em,topsep=.4em]\item[A.] A study of the fast-food service.
\item[B.] Fast food customer satisfaction.
\item[C.] McDonald's new business strategies.
\item[D.] Competition in the fast-food industry.
\end{itemize}

4.


\begin{itemize}[leftmargin=2em,itemsep=.2em,topsep=.4em]\item[A.] Customers' higher demands.
\item[B.] The inefficiency of employees.
\item[C.] Increased variety of products.
\item[D.] The rising number of customers.
\end{itemize}

\textbf{Questions 5 to 7 are based on the news report you have just heard.}


5.


\begin{itemize}[leftmargin=2em,itemsep=.2em,topsep=.4em]\item[A.] International treaties regarding space travel programs.
\item[B.] Legal issues involved in commercial space exploration.
\item[C.] U.S. government's approval of private space missions.
\item[D.] Competition among public and private space companies.
\end{itemize}

6.


\begin{itemize}[leftmargin=2em,itemsep=.2em,topsep=.4em]\item[A.] Deliver scientific equipment to the moon.
\item[B.] Approve a new mission to travel into outer space.
\end{itemize}

1


\begin{itemize}[leftmargin=2em,itemsep=.2em,topsep=.4em]\item[C.] Work with federal agencies on space programs.
\item[D.] Launch a manned spacecraft to Mars.
\end{itemize}

\noindent\begin{tabularx}{\linewidth}{@{}lXXXX@{}}7. & \textbf{A.} It is significant. & \textbf{B.} It is promising. & \textbf{C.} It is unpredictable. & \textbf{D.} It is unprofitable.\\\end{tabularx}\par

\subsection*{\textbf{Section B}}


\textbf{Directions:}\textit{ In this section, you will hear two long conversations. At the end of each} \textit{conversation, you will hear four questions. Both the conversation and the questions will bespoken} \textit{only once. After you hear a question, you must choose the best answer from the four} \textit{choices marked A), B), C) and D). Then mark the corresponding letter on Answer Sheet 1with a si} \textit{ngle line through the centre.}


\textbf{Questions 8 to 11 are based on the conversation you have just heard.}


8.


\begin{itemize}[leftmargin=2em,itemsep=.2em,topsep=.4em]\item[A.] Visiting her family in Thailand.
\item[B.] Showing friends around Phuket.
\item[C.] Swimming around a Thai island.
\item[D.] Lying in the sun on a Thai beach.
\end{itemize}

9.


\begin{itemize}[leftmargin=2em,itemsep=.2em,topsep=.4em]\item[A.] She visited a Thai orphanage.
\item[B.] She met a Thai girl's parents.
\item[C.] She learned some Thai words.
\item[D.] She sunbathed on a Thai beach.
\end{itemize}

10.


\begin{itemize}[leftmargin=2em,itemsep=.2em,topsep=.4em]\item[A.] His class will start in a minute.
\item[B.] He has got an incoming phone call.
\item[C.] Someone is knocking at his door.
\item[D.] His phone is running out of power.
\end{itemize}

11.


\begin{itemize}[leftmargin=2em,itemsep=.2em,topsep=.4em]\item[A.] He is interested in Thai artworks.
\item[B.] He is going to open a souvenir shop.
\item[C.] He collects things from different countries.
\item[D.] He wants to know more about Thai culture.
\end{itemize}

\textbf{Questions 12 to 15 are based on the conversation you have just heard.}


12.


\begin{itemize}[leftmargin=2em,itemsep=.2em,topsep=.4em]\item[A.] Buying some fitness equipment for the new gym.
\item[B.] Opening a gym and becoming personal trainers.
\item[C.] Signing up for a weight-loss course.
\item[D.] Trying out a new gym in town.
\end{itemize}

13.


\begin{itemize}[leftmargin=2em,itemsep=.2em,topsep=.4em]\item[A.] Professional personal training.
\item[B.] Free exercise for the first week.
\item[C.] A discount for a half-year membership.
\item[D.] Additional benefits for young couples.
\end{itemize}

14.


\begin{itemize}[leftmargin=2em,itemsep=.2em,topsep=.4em]\item[A.] The safety of weight-lifting.
\item[B.] The high membership fee.
\item[C.] The renewal of his membership.
\item[D.] The operation of fitness equipment.
\end{itemize}

15.


\begin{itemize}[leftmargin=2em,itemsep=.2em,topsep=.4em]\item[A.] She wants her invitation renewed.
\item[B.] She used to do 200 sit-ups every day.
\end{itemize}

2


\begin{itemize}[leftmargin=2em,itemsep=.2em,topsep=.4em]\item[C.] She knows the basics of weight-lifting.
\item[D.] She used to be the gym's personal trainer.
\end{itemize}

\subsection*{\textbf{Section C}}


\textbf{Directions:}\textit{ In this section, you will hear three passages. At the end of each passage, you will} \textit{hear three or four questions. Both the passage and the questions will be spoken only once.After yo} \textit{u hear a question, you must choose the best answer from the four choices marked A),B), C) and D).} \textit{Then mark the corresponding letter on Answer Sheet 1 with a single line through the centre.}


\textbf{Questions 16 to 18 are based on the passage you have just heard.}


16.


\begin{itemize}[leftmargin=2em,itemsep=.2em,topsep=.4em]\item[A.] They tend to be nervous during interviews.
\item[B.] They often apply for a number of positions.
\item[C.] They worry about the results of their applications.
\item[D.] They search extensively for employers' information.
\end{itemize}

17.


\begin{itemize}[leftmargin=2em,itemsep=.2em,topsep=.4em]\item[A.] Get better organized.
\item[B.] Edit their references.
\item[C.] Find better-paid jobs.
\item[D.] Analyze the searching process.
\end{itemize}

18.


\begin{itemize}[leftmargin=2em,itemsep=.2em,topsep=.4em]\item[A.] Provide their data in detail.
\item[B.] Personalize each application.
\item[C.] Make use of better search engines.
\item[D.] Apply for more promising positions.
\end{itemize}

\textbf{Questions 19 to 21 are based on the passage you have just heard.}


19.


\begin{itemize}[leftmargin=2em,itemsep=.2em,topsep=.4em]\item[A.] If kids did not like school, real learning would not take place.
\item[B.] If not forced to go to school, kids would be out in the streets.
\item[C.] If schools stayed the way they are, parents were sure to protest.
\item[D.] If teaching failed to improve, kids would stay away from school.
\end{itemize}

20.


\begin{itemize}[leftmargin=2em,itemsep=.2em,topsep=.4em]\item[A.] Allow them to play interesting games in class.
\item[B.] Try to stir up their interest in lab experiments.
\item[C.] Let them stay home and learn from their parents.
\item[D.] Design activities they now enjoy doing on holidays.
\end{itemize}

21.


\begin{itemize}[leftmargin=2em,itemsep=.2em,topsep=.4em]\item[A.] Allow kids to learn at their own pace.
\item[B.] Encourage kids to learn from each other.
\item[C.] Organize kids into various interest groups.
\item[D.] Take kids out of school to learn at first hand.
\end{itemize}

\textbf{Questions 22 to 25 are based on the passage you have just heard.}


22.


\begin{itemize}[leftmargin=2em,itemsep=.2em,topsep=.4em]\item[A.] It is especially popular in Florida and Alaska.
\item[B.] It is a major social activity among the young.
\item[C.] It is seen almost anywhere and on any occasion.
\item[D.] It is even more expressive than the written word.
\end{itemize}

23.


\begin{itemize}[leftmargin=2em,itemsep=.2em,topsep=.4em]\item[A.] It is located in a big city in Iowa.
\item[B.] It is really marvelous to look at.
\item[C.] It offers free dance classes to seniors.
\item[D.] It offers people a chance to socialize.
\end{itemize}

24.


\begin{itemize}[leftmargin=2em,itemsep=.2em,topsep=.4em]\item[A.] Their state of mind improved.
\item[B.] They became better dancers.
\end{itemize}

3


\begin{itemize}[leftmargin=2em,itemsep=.2em,topsep=.4em]\item[C.] They enjoyed better health.
\item[D.] Their relationship strengthened.
\end{itemize}

\noindent\begin{tabularx}{\linewidth}{@{}lXXXX@{}}25. & \textbf{A.} It is fun. & \textbf{B.} It is life. & \textbf{C.} It is exhausting. & \textbf{D.} It is rhythmical.\\\end{tabularx}\par

\subsection*{\textbf{Part} III \textbf{Reading Comprehension} \textbf{(40 minutes)}}


\subsection*{\textbf{Section A}}


\textbf{Directions:}\textit{ In this section, there is a passage with10 blanks. You are required to select one word}


\textit{for each blank from a list of choices given in a word bank following the passage. Read the}


\subsection*{\textit{passage through carefully before making your choices. Each choice in the bank is identified}}


\textit{by a letter. Please mark the corresponding letter for each item on Answer Sheet 2 with a}


\textit{single line through the centre. You may not use any of the words in the bank more than once.}


Since the 1940s, southern California has had a reputation for smog. Things are not as bad as they once were but, according to the American Lung Association, Los Angeles is still the worst city in the United States for levels of \_\_26\_\_ . Gazing down on the city from the Getty Center, an art museum in the Santa Monica Mountains, one would find the view of the Pacific Ocean blurred by the haze (霾). Nor is the state's bad air \_\_27\_\_ to its south. Fresno, in the central valley, comes top of the list in America for year-round pollution. Residents' hearts and lungs are affected as a \_\_28\_\_ . All of which, combined with California's reputation as the home of technological \_\_29\_\_ , makes the place ideal for developing and testing systems designed to monitor pollution in \_\_30\_\_ . And that is just what Aclima, a new firm in San Francisco, has been doing over the past few months. It has been trying out monitoring stations that are \_\_31\_\_ to yield minute-to-minute maps of \_\_32\_\_ air pollution. Such stations will also be able to keep an eye on what is happening inside buildings, including offices.


To this end, Aclima has been \_\_33\_\_ with Google's Street View system. Davida Herzl, Aclima's boss, says they have revealed pollution highs on days when San Francisco's transit workers went on strike and the city's \_\_34\_\_ were forced to use their cars. Conversely, "cycle to work" days have done their job by \_\_35\_\_ pollution lows.


\begin{itemize}[leftmargin=2em,itemsep=.2em,topsep=.4em]\item[A.] assisted
\item[B.] collaborating
\item[C.] consequence
\item[D.] consumers
\item[E.] creating
\item[F.] detail
\item[G.] domestic
\item[H.] frequently
\item[I.] inhabitants
\item[J.] innovation
\item[K.] intended
\item[L.] outdoor
\item[M.] pollutants
\item[N.] restricted
\item[O.] sum
\end{itemize}

\subsection*{\textbf{Section B}}


\textbf{Directions:}\textit{ In this section, you are going to read a passage with ten statements attached to} \textit{it. Each statement contains information given in one of the paragraphs. Identify the} \textit{paragraph from which the information is derived. You may choose a paragraph more than} \textit{once. Each paragraph is marked with a letter. Answer the questions by marking the} \textit{corresponding letter on Answer Sheet 2.}


\textbf{As Tourists Crowd Out Locals, Venice Faces 'Endangered' List}


\begin{itemize}[leftmargin=2em,itemsep=.2em,topsep=.4em]\item[A.] On a recent fall morning, a large crowd blocked the steps at one of Venice's main tourist
\end{itemize}

sites, the Rialto Bridge. The Rialto Bridge is one of the four bridges spanning the Grand Canal. It is the oldest bridge across the canal, and was the dividing line between the districts of San Marco and San Polo. But on this day, there was a twist: it was filled with Venetians, not tourists.


4


\begin{itemize}[leftmargin=2em,itemsep=.2em,topsep=.4em]\item[B.] "People are cheering and holding their carts in the air," says Giovanni Giorgio, who
\end{itemize}

helped organize the march with a grass-roots organization called Generazione '90. The carts he refers to are small shopping carts—the symbol of a true Venetian. "It started as a joke," he says with a laugh. "The idea was to put blades on the wheels! You know? Like Ben Hur. Precisely like that, you just go around and run people down."


\begin{itemize}[leftmargin=2em,itemsep=.2em,topsep=.4em]\item[C.] Venice is one of the hottest tourist destinations in the world. But that's a problem. Up to
\end{itemize}

90,000 tourists crowd its streets and canals every day—far outnumbering the 55,000 permanent residents. The tourist increase is one key reason the city's population is down from 175,000 in the 1950s. The outnumbered Venetians have been steadily fleeing. And those who stick around are tired of living in a place where they can't even get to the market without swimming through a sea of picture-snapping tourists. Imagine, navigating through 50,000 people while on the way to school or to work.


\begin{itemize}[leftmargin=2em,itemsep=.2em,topsep=.4em]\item[D.] Laura Chigi, a grandmother at the march, says the local and national governments have
\end{itemize}

failed to do anything about the crowds for decades, because they're only interested in tourism—the primary industry in Venice, worth more than \$3 billion in 2015. "Venice is a cash cow," she says, "and everyone wants a piece."


\begin{itemize}[leftmargin=2em,itemsep=.2em,topsep=.4em]\item[E.] Just beyond St. Mark's Square, a cruise ship passes, one of hundreds every year that
\end{itemize}

appear over their medieval (中世纪的) surroundings. Their massive wake creates waves at the bottom of the sea, weakening the foundations of the centuries-old buildings themselves. "Every time I see a cruise ship, I feel sad," Chigi says. "You see the mud it drags; the destruction it leaves in its wake? That hurts the ancient wooden poles holding up the city underwater. One day we'll see Venice break down."


\begin{itemize}[leftmargin=2em,itemsep=.2em,topsep=.4em]\item[F.] For a time, UNESCO, the cultural wing of the United Nations, seemed to agree. Two years
\end{itemize}

ago, it put Italy on notice, saying the government was not protecting Venice. UNESCO considers the entire city a World Heritage Site, a great honor that means Venice, at the cultural level, belongs to all of the world's people. In 2014, UNESCO gave Italy two years to manage Venice's flourishing tourism or the city would be placed on another list—World Heritage In Danger, joining such sites as Aleppo and Palmyra, destroyed by the war in Syria.


\begin{itemize}[leftmargin=2em,itemsep=.2em,topsep=.4em]\item[G.] Venice's deadline passed with barely a murmur (嘟哝) this summer, just as UNESCO was
\end{itemize}

meeting in Istanbul. Only one representative, Jad Tabet from Lebanon, tried to raise the issue. "For several years, the situation of heritage in Venice has been worsening, and it has now reached a dramatic situation," Tabet told UNESCO. "We have to act quickly—there is not a moment to waste."


\begin{itemize}[leftmargin=2em,itemsep=.2em,topsep=.4em]\item[H.] But UNESCO didn't even hold a vote. "It's been postponed until 2017," says Anna Somers,
\end{itemize}

the founder and CEO of The Art Newspaper and the former head of Venice in Peril, a group devoted to restoring Venetian art. She says the main reason the U.N. cultural organization didn't vote to declare Venice a World Heritage Site In Danger is because UNESCO has become


5


"intensely politicized. There would have been some back-room negotiations."


\begin{itemize}[leftmargin=2em,itemsep=.2em,topsep=.4em]\item[I.] Italy boasts more UNESCO World Heritage Sites than any other country in the world,
\end{itemize}

granting it considerable power and influence within the organization. The former head of the UNESCO World Heritage Centre, which oversees heritage sites, is Francesco Bandarin, a Venetian who now serves as UNESCO's assistant director-general for culture.


\begin{itemize}[leftmargin=2em,itemsep=.2em,topsep=.4em]\item[J.] Earlier this year, Italy signed an accord with UNESCO to establish a task force of police
\end{itemize}

art detectives and archaeologists (考古学家) to protect cultural heritage from natural disasters and terror groups, such as ISIS. The accord underlined Italy's global reputation as a good steward of art and culture.


\begin{itemize}[leftmargin=2em,itemsep=.2em,topsep=.4em]\item[K.] But adding Venice to the UNESCO endangered list—which is dominated by sites in
\end{itemize}

developing and conflict-ridden countries—would be an international embarrassment, and could even hurt Italy's profitable tourism industry. The Italian Culture Ministry says it is unaware of any government efforts to pressure UNESCO. As for the organization itself, it declined a request for an interview.


\begin{itemize}[leftmargin=2em,itemsep=.2em,topsep=.4em]\item[L.] The city's current mayor, Luigi Brugnaro, has ridiculed UNESCO and told it to mind its
\end{itemize}

own business, while continuing to support the cruise ship industry, which employs 5,000 Venice residents.


\begin{itemize}[leftmargin=2em,itemsep=.2em,topsep=.4em]\item[M.] As for Venetians, they're beyond frustrated and hoping for a solution soon. "It's a
\end{itemize}

nightmare for me. Some situations are really difficult with tourists around," says Giorgio as he navigates around a swelling crowd at the Rialto Bridge. "There are just so many of them. They never know where they are going, and do not walk in an orderly manner. Navigating the streets can be exhausting."


\begin{itemize}[leftmargin=2em,itemsep=.2em,topsep=.4em]\item[N.] Then it hits him: This crowd isn't made up of tourists. They're Venetians. Giorgio says he's
\end{itemize}

never experienced the Rialto Bridge this way in all his 22 years. "For once, we are the ones who are blocking the traffic," he says delightedly. "It feels unreal. It feels like we're some form of endangered species. It's just nice. The feeling is just pure." But, he worries, if tourism isn't managed and his fellow locals continue to move to the mainland, his generation might be the last who can call themselves native Venetians.


36. The passing cruise ships will undermine the foundations of the ancient buildings in


Venice.


37. The Italian government has just reached an agreement with UNESCO to take measures to


protect its cultural heritage.


38. The heritage situation in Venice has been deteriorating in the past few years.


39. The decrease in the number of permanent residents in Venice is mainly due to the increase


of tourists.


40. If tourism gets out of control, native Venetians may desert the city altogether one day.


41. UNESCO urged the Italian government to undertake its responsibility to protect Venice.


6


42. The participants in the Venetian march used shopping carts to show they were 100\% local


residents.


43. Ignoring UNESCO's warning, the mayor of Venice maintains his support of the city's


tourism industry.


44. One woman says that for decades the Italian government and local authorities have only


focused on the revenues from tourism.


45. UNESCO has not yet decided to put Venice on the list of World Heritage Sites In Danger.


\subsection*{\textbf{Section C}}


\textbf{Directions:}\textit{ There are 2 passages in this section. Each passage is followed by some questions} \textit{or unfinished statements. For each of them there are four choices marked A), B), C) and}


\begin{itemize}[leftmargin=2em,itemsep=.2em,topsep=.4em]\item[D.] \textit{. You should decide on the best choice and mark the corresponding letter on Answer Sheet2 wit} \textit{h a single line through the centre.}
\end{itemize}

\subsection*{\textbf{Passage One}}


\textbf{Questions 46 and 50 are based on the following passage.}


\subsection*{Passage One}


Questions 46 to 50 are based on the following passage.


Living in an urban area with green spaces has a long-lasting positive impact on people's mental well-being, a study has suggested. UK researchers found moving to a green space had a sustained positive effect, unlike pay rises or promotions, which only provided a short-term boost. Co-author Mathew White, from the University of Exeter, UK, explained that the study showed people living in greener urban areas were displaying fewer signs of depression or anxiety. "There could be a number of reasons," he said, "for example, people do many things to make themselves happier: they strive for promotion or pay rises, or they get married. But the trouble with those things is that within six months to a year, people are back to their original baseline levels of well-being. So, these things are not sustainable; they don't make us happy in the long term. We found that for some lottery (彩票) winners who had won more than ￡500,000 the positive effect was definitely there, but after six months to a year, they were back to the baseline."


Dr. White said his team wanted to see whether living in greener urban areas had a lasting positive effect on people's sense of well-being or whether the effect also disappeared after a period of time. To do this, the team used data from the British Household Panel Survey compiled by the University of Essex.


Explaining what the data revealed, he said: "What you see is that even after three years, mental health is still better, which is unlike many other things that we think will make us happy." He observed that people living in green spaces were less stressed, and less stressed people made more sensible decisions and communicated better.


With a growing body of evidence establishing a link between urban green spaces and a positive impact on human well-being, Dr. White said, "There's growing interest among public policy officials, but the trouble is who funds it. What we really need at a policy level is to decide where the money will come from to help support good quality local green spaces."


7


46.According to one study, what do green spaces do to people?


\begin{itemize}[leftmargin=2em,itemsep=.2em,topsep=.4em]\item[A.] Improve their work efficiency.
\item[B.] Add to their sustained happiness.
\item[C.] Help them build a positive attitude towards life.
\item[D.] Lessen their concerns about material well-being.
\end{itemize}

47.What does Dr. White say people usually do to make themselves happier?


\begin{itemize}[leftmargin=2em,itemsep=.2em,topsep=.4em]\item[A.] Earn more money.
\item[B.] Settle in an urban area.
\item[C.] Gain fame and popularity.
\item[D.] Live in a green environment.
\end{itemize}

48.What does Dr. White try to find out about living in a greener urban area?


\begin{itemize}[leftmargin=2em,itemsep=.2em,topsep=.4em]\item[A.] How it affects different people.
\item[B.] How strong its positive effect is.
\item[C.] How long its positive effect lasts.
\item[D.] How it benefits people physically.
\end{itemize}

49.What did Dr. White's research reveal about people living in a green environment?


\begin{itemize}[leftmargin=2em,itemsep=.2em,topsep=.4em]\item[A.] Their stress was more apparent than real.
\item[B.] Their decisions required less deliberation.
\item[C.] Their memories were greatly strengthened.
\item[D.] Their communication with others improved.
\end{itemize}

50.According to Dr. White, what should the government do to build more green spaces in cities?


\begin{itemize}[leftmargin=2em,itemsep=.2em,topsep=.4em]\item[A.] Find financial support.
\item[B.] Improve urban planning.
\item[C.] Involve local residents in the effort.
\item[D.] Raise public awareness of the issue.Passage Two
\end{itemize}

\subsection*{\textbf{Passage Two}}


\textbf{Questions 51 to 55 are based on the following passage.}


You probably know about the Titanic, but it was actually just one of three state-of-the-art (先进的) ocean ships back in the day. The Olympic class ships were built by the Harland \&amp; Wolff ship makers in Northern Ireland for the White Star Line company. The Olympic class included the Olympic, the Britannic and the Titanic. What you may not know is that the Titanic wasn't even the flagship of this class. All in all, the Olympic class ships were marvels of sea engineering, but they seemed cursed to suffer disastrous fates.


The Olympic launched first in 1910, followed by the Titanic in 1911, and lastly the Britannic in 1914. The ships had nine decks, and White Star Line decided to focus on making them the most luxurious ships on the water.


Stretching 269.13 meters, the Olympic class ships were wonders of naval technology, and everyone thought that they would continue to be so for quite some time. However, all suffered


8


terrible accidents on the open seas. The Olympic got wrecked before the Titanic did, but it was the only one to survive and maintain a successful career of 24 years. The Titanic was the first to sink after famously hitting a huge iceberg in 1912. Following this disaster, the Britannic hit a naval mine in 1916 and subsequently sank as well.


Each ship was coal-powered by several boilers constantly kept running by exhausted crews below deck. Most recognizable of the ship designs are the ship's smoke stacks, but the fourth stack was actually just artistic in nature and served no functional purpose. While two of these ships sank, they were all designed with double hulls (船体) believed to make them "unsinkable", perhaps a mistaken idea that led to the Titanic's and the Britannic's tragic end.


The Olympic suffered two crashes with other ships and went on to serve as a hospital ship and troop transport in World War I. Eventually, she was taken out of service in 1935, ending the era of the luxurious Olympic class ocean liners.


51.What does the passage say about the three Olympic class ships?


\begin{itemize}[leftmargin=2em,itemsep=.2em,topsep=.4em]\item[A.] They performed marvellously on the sea.
\item[B.] They could all break the ice in their way.
\item[C.] They all experienced terrible misfortunes.
\item[D.] They were models of modern engineering.
\end{itemize}

52.What did White Star Line have in mind when it purchased the three ships?


\begin{itemize}[leftmargin=2em,itemsep=.2em,topsep=.4em]\item[A.] Their capacity of sailing across all waters.
\item[B.] The utmost comfort passengers could enjoy.
\item[C.] Their ability to survive disasters of any kind.
\item[D.] The long voyages they were able to undertake.
\end{itemize}

53.What is said about the fourth stack of the ships?


\begin{itemize}[leftmargin=2em,itemsep=.2em,topsep=.4em]\item[A.] It was a mere piece of decoration.
\item[B.] It was the work of a famous artist.
\item[C.] It was designed to let out extra smoke.
\item[D.] It was easily identifiable from afar.
\end{itemize}

54.What might have led to the tragic end of the Titanic and the Britannic?


\begin{itemize}[leftmargin=2em,itemsep=.2em,topsep=.4em]\item[A.] Their unscientific designs.
\item[B.] Their captains' misjudgment.
\item[C.] The assumption that they were built with the latest technology.
\item[D.] The belief that they could never sink with a double-layer body.
\end{itemize}

55.What happened to the ship Olympic in the end?


\begin{itemize}[leftmargin=2em,itemsep=.2em,topsep=.4em]\item[A.] She was used to carry troops.
\item[B.] She was sunk in World War
\item[C.] She was converted into a hospital ship.
\item[D.] She was retired after her naval service.
\item[I.] 
\end{itemize}

9


\subsection*{\textbf{Part IV} \textbf{Translation} \textbf{(30 minutes)}}


\textbf{Directions:}\textit{ For this part, you are allowed 30minutes to translate a passage from Chinese} \textit{into} \textit{English. You should write your answer on Answer Sheet 2.}


公交车曾是中国人出行的主要交通工具。近年来，由于私家车数量不断增多，城市的交通问题越来越严重。许多城市为了鼓励更多人乘坐公交车出行，一直在努力改善公交车的服务质量。车辆的设施不断更新，车速也有了显著提高。然而，公交车的票价却依然相当低廉。 现在，在大多数城市，许多当地老年市民都可以免费乘坐公交车。


10

\end{document}